\section{Related Work}
\label{sec:related_work}

Automated traffic surveillance through computer vision spans three closely intertwined technical domains: deep real-time object detection, multi-object tracking (MOT) data association, and virtual counting line geometries.

\subsection{Real-Time Deep Object Detection}
Object detection in intelligent transportation systems has evolved rapidly from classical sliding-window feature descriptors toward end-to-end deep convolutional neural networks. The You Only Look Once (YOLO) family of single-stage detectors, introduced by Redmon et al.~\cite{redmon2016yolo}, framed object localization and classification as a unified bounding box regression task. The architecture was subsequently refined through YOLOv3~\cite{redmon2018yolov3}, which introduced multi-scale feature pyramids. Later iterations, including YOLOv4~\cite{bochkovskiy2020yolov4} and YOLOv7~\cite{wang2023yolov7}, incorporated path aggregation networks, cross-stage partial connections, and advanced data augmentation strategies to further optimize the Pareto frontier between detection precision and latency. Modern implementations trained on benchmark corpora such as Microsoft COCO~\cite{lin2014coco} deliver high mean Average Precision (mAP) across canonical vehicle categories, including cars, motorcycles, buses, and trucks. 

In the baseline study by Majumder and Wilmot~\cite{majumder2023automated}, a Darknet implementation of YOLOv3 converted to TensorFlow served as the primary detector. However, their model suffered from substantial computational overhead, requiring 1.7 hours to process a single hour of footage on a contemporary consumer GPU, and treated all vehicles as an undifferentiated single class. In contrast, our approach utilizes modern, lightweight YOLO architectures supporting native multi-class localization and sub-10 millisecond inference latencies on commodity GPU hardware, enabling real-time streaming operations without offline pre-processing bottlenecks.

\subsection{Multi-Object Tracking (MOT) in Roadway Scenes}
The tracking-by-detection paradigm dominates multi-target video analysis. Bewley et al.~\cite{bewley2016sort} established Simple Online and Realtime Tracking (SORT), combining a discrete linear Kalman filter~\cite{kalman1960new} for state propagation with the Hungarian algorithm~\cite{kuhn1955hungarian} for spatial bounding-box association via Intersection over Union (IoU). While exceptionally fast, SORT exhibits severe trajectory fragmentation when targets undergo visual occlusion, as bounding boxes failing to meet high detection thresholds are discarded entirely. Wojke et al.~\cite{wojke2017deepsort} introduced DeepSORT, incorporating deep appearance feature embeddings to mitigate identity switches across prolonged occlusions. However, extracting deep re-identification (ReID) feature vectors for every candidate box imposes non-trivial computational overhead during edge execution.

Recently, Zhang et al.~\cite{zhang2022bytetrack} introduced ByteTrack, challenging the assumption that low-score detection boxes represent background noise. ByteTrack adopts a hierarchical, two-stage association strategy: first matching high-confidence detections to active tracklets, and subsequently matching unmatched tracklets with low-confidence detections. This simple yet principled mechanism recovers partially occluded objects without invoking computationally intensive ReID neural networks. Aharon et al.~\cite{aharon2022botsort} extended this concept with BoT-SORT, integrating Global Motion Compensation (GMC) via camera feature-point tracking. While advantageous for moving or hand-held cameras, our analysis demonstrates that GMC incurs significant and unnecessary latency when applied to fixed roadside surveillance cameras. Consequently, our framework adopts ByteTrack as the primary tracking engine to ensure occlusion resilience while preserving maximal pipeline throughput.

\subsection{Virtual Tripwires and Traffic Flow Estimation}
Virtual Detection Lines (VDLs) and virtual loop detectors have long served as non-intrusive computational analogues to physical inductive loops. Early systems sampled 1D pixel profiles over time to generate Time-Spatial Images (TSIs), detecting vehicle transits through spatio-temporal slice correlations. While lightweight, TSI techniques degrade under multi-lane occlusion, vehicle shadow interference, and lateral lane drift.

Modern tripwire systems map tracked vehicle trajectories across virtual polyline gates defined within the image coordinate frame. The baseline formulation introduced by Majumder and Wilmot~\cite{majumder2023automated} implemented a three-line virtual gate comprising a user-drawn mid-line flanked by two automatically projected parallel boundary lines. Directional transit was determined by evaluating the order in which the vehicle's bounding box centroid intersected the boundary line and the mid-line. While conceptually sound, their implementation suffered from systematic undercounting arising from three distinct mechanisms: (i) bounding box centroids introduced severe perspective-induced timing errors for high-profile commercial vehicles, (ii) tracking failures under vehicle overlap led to dropped crossing events, and (iii) vehicles that entered the detection frame between the outer gate and the mid-line were permanently ignored. 

Our framework directly addresses and rectifies each of these systemic failure modes. By tracking the tire-to-pavement contact patch ($x_{\text{mid}}, y_{\text{max}}$), computing signed 2D cross-product distances, providing adaptive boundary fallback logic, and decoupling interval aggregation from host clock drift, \texttt{vcount} provides a mathematically robust and reproducible counting architecture.
